\documentclass[11pt,a4paper]{article}

\usepackage[margin=1in]{geometry}
\usepackage{graphicx}
\usepackage{listings}
\usepackage{xcolor}
\usepackage{booktabs}
\usepackage{longtable}
\usepackage{hyperref}
\usepackage{parskip}
\usepackage{fancyhdr}
\usepackage{amsmath}

\definecolor{codebg}{RGB}{245,245,245}
\definecolor{codekw}{RGB}{0,0,180}
\definecolor{codecomment}{RGB}{100,100,100}
\definecolor{codestr}{RGB}{163,21,21}

\lstset{
  backgroundcolor=\color{codebg},
  basicstyle=\ttfamily\small,
  keywordstyle=\color{codekw}\bfseries,
  commentstyle=\color{codecomment}\itshape,
  stringstyle=\color{codestr},
  breaklines=true,
  frame=single,
  language=Python,
  numbers=left,
  numberstyle=\tiny,
  xleftmargin=2em,
  showstringspaces=false
}

\pagestyle{fancy}
\fancyhf{}
\rhead{ICS1512 --- ML Algorithms Lab}
\lhead{Experiment 5}
\cfoot{\thepage}

\title{\textbf{Experiment 5: Decision Tree and Random Forest --- \\ A Comparative Classification Study}}
\author{
  M.Tech (Integrated) Computer Science \& Engineering \\
  ICS1512 --- Machine Learning Algorithms Laboratory \\
  Sri Sivasubramaniya Nadar College of Engineering, Chennai
}
\date{Academic Year 2026--2027 (Odd) \quad Batch: 2024--2029}

\begin{document}
\maketitle

\section*{Objective}
\begin{itemize}
  \item Implement a Decision Tree classifier.
  \item Extend the Decision Tree into a Random Forest ensemble model.
  \item Study the impact of hyperparameters on overfitting and generalization.
  \item Select optimal hyperparameters using 5-Fold Cross-Validation.
  \item Compare single-tree and ensemble-tree models.
\end{itemize}

\section*{Dataset}
Wisconsin Diagnostic Breast Cancer Dataset --- 569 samples, 30 numerical features, target classes
Malignant (M) and Benign (B). Loaded via \texttt{sklearn.datasets.load\_breast\_cancer}, which provides
the exact UCI dataset (\url{https://archive.ics.uci.edu}) bundled directly with scikit-learn.

\section{Theory}

\subsection{Decision Tree Classifier}
A Decision Tree recursively splits the feature space using impurity measures (Gini index or entropy)
to form a set of decision rules. Deeper trees can fit the training data almost perfectly but are prone
to \textbf{high variance and overfitting}, since they end up modeling noise in addition to genuine
patterns.

\subsection{Random Forest Classifier}
Random Forest is an ensemble technique that trains many decision trees on bootstrapped samples of the
training data (\textbf{bagging}) and considers a random subset of features at each split
(\textbf{random feature selection}). Predictions are combined by \textbf{majority voting}. This reduces
variance and improves generalization relative to a single decision tree.

\section{Methodology}
\begin{enumerate}
  \item Loaded the dataset and encoded class labels (0 = Malignant, 1 = Benign).
  \item Performed EDA: class distribution and feature correlation analysis.
  \item Split the dataset into training and testing sets (80--20, stratified).
  \item Trained a baseline Decision Tree classifier.
  \item Defined a hyperparameter search space for the Decision Tree.
  \item Performed 5-Fold Cross-Validation (via \texttt{GridSearchCV}) to evaluate each combination.
  \item Selected the hyperparameters yielding the best average CV accuracy.
  \item Trained a Random Forest classifier (baseline and hyperparameter-tuned).
  \item Repeated cross-validation-based hyperparameter selection for the Random Forest.
  \item Compared both tuned models using accuracy, precision, recall, F1-score, confusion matrix, and
        ROC/AUC.
\end{enumerate}

\section{Exploratory Data Analysis}

The dataset is moderately imbalanced towards the benign class. Figure~\ref{fig:class_dist} shows the
class distribution, and Figure~\ref{fig:corr} shows pairwise correlations among the first 15 features
(several \texttt{mean}, \texttt{worst}, and \texttt{error} feature groups are highly correlated with
each other, as expected). Figure~\ref{fig:boxplots} shows that \texttt{mean radius} and \texttt{mean
concavity} separate the two classes noticeably better than \texttt{mean texture}.

\begin{figure}[h]
  \centering
  \includegraphics[width=0.55\linewidth]{figs/fig_01.png}
  \caption{Class distribution: Malignant vs.\ Benign.}
  \label{fig:class_dist}
\end{figure}

\begin{figure}[h]
  \centering
  \includegraphics[width=0.85\linewidth]{figs/fig_02.png}
  \caption{Feature correlation heatmap (first 15 features).}
  \label{fig:corr}
\end{figure}

\begin{figure}[h]
  \centering
  \includegraphics[width=0.95\linewidth]{figs/fig_03.png}
  \caption{Distribution of mean radius, mean texture, and mean concavity by diagnosis.}
  \label{fig:boxplots}
\end{figure}

\clearpage
\section{Baseline Models}

An unconstrained baseline Decision Tree reached \textbf{100\% training accuracy} but only
\textbf{91.23\% test accuracy} (depth 7, 19 leaves) --- a clear overfitting signature. The baseline
Random Forest (default settings) reached 100\% training accuracy and \textbf{95.61\% test accuracy},
already showing a smaller generalization gap purely from ensembling. Figure~\ref{fig:tree} shows the
top two levels of the baseline Decision Tree.

\begin{figure}[h]
  \centering
  \includegraphics[width=0.95\linewidth]{figs/fig_04.png}
  \caption{Baseline Decision Tree — top two levels.}
  \label{fig:tree}
\end{figure}

\section{Hyperparameter Tuning Results}

\subsection{Decision Tree Cross-Fold Results}
Search space: \texttt{criterion} $\in \{$gini, entropy$\}$, \texttt{max\_depth} $\in \{3,5,7,10,$None$\}$,
\texttt{min\_samples\_split} $\in \{2,5,10\}$, \texttt{min\_samples\_leaf} $\in \{1,2,4\}$ ---
90 combinations evaluated via 5-fold cross-validation. Table~\ref{tab:dt_cv} shows the best-scoring
combination for each (criterion, max\_depth) pair.

\begin{table}[h]
\centering
\caption{Decision Tree Hyperparameter Evaluation using 5-Fold Cross-Validation (top results)}
\label{tab:dt_cv}
\begin{tabular}{llcc}
\toprule
\textbf{Criterion} & \textbf{Max Depth} & \textbf{Avg CV Accuracy (\%)} & \textbf{Avg CV F1 Score} \\
\midrule
gini    & 10        & 93.85 & 0.9506 \\
gini    & 5         & 93.85 & 0.9506 \\
gini    & 7         & 93.85 & 0.9506 \\
gini    & None      & 93.85 & 0.9506 \\
entropy & 3         & 93.41 & 0.9480 \\
gini    & 3         & 93.19 & 0.9455 \\
entropy & None      & 92.97 & 0.9435 \\
entropy & 5         & 92.97 & 0.9439 \\
\bottomrule
\end{tabular}
\end{table}

\textbf{Best Decision Tree hyperparameters:} \texttt{criterion = gini}, \texttt{max\_depth = 5},
\texttt{min\_samples\_leaf = 4}, \texttt{min\_samples\_split = 2}, with best average CV accuracy of
\textbf{93.85\%}.

\subsection{Random Forest Cross-Fold Results}
Search space: \texttt{n\_estimators} $\in \{50,100,200\}$, \texttt{max\_depth} $\in \{5,10,$None$\}$,
\texttt{max\_features} $\in \{$sqrt, log2$\}$, \texttt{bootstrap} $\in \{$True, False$\}$ ---
36 combinations evaluated via 5-fold cross-validation. Table~\ref{tab:rf_cv} shows the top results.

\begin{table}[h]
\centering
\caption{Random Forest Hyperparameter Evaluation using 5-Fold Cross-Validation (top results)}
\label{tab:rf_cv}
\begin{tabular}{lllcc}
\toprule
\textbf{n\_estimators} & \textbf{Max Depth} & \textbf{Max Features} & \textbf{Avg CV Accuracy (\%)} & \textbf{Avg CV F1 Score} \\
\midrule
100 & 10   & log2 & 96.48 & 0.9718 \\
100 & None & log2 & 96.48 & 0.9718 \\
200 & 10   & log2 & 96.26 & 0.9703 \\
50  & None & log2 & 96.26 & 0.9702 \\
100 & None & sqrt & 96.26 & 0.9699 \\
50  & 10   & log2 & 96.26 & 0.9702 \\
200 & None & sqrt & 96.26 & 0.9700 \\
100 & 10   & sqrt & 96.26 & 0.9699 \\
\bottomrule
\end{tabular}
\end{table}

\textbf{Best Random Forest hyperparameters:} \texttt{n\_estimators = 100}, \texttt{max\_depth = 10},
\texttt{max\_features = log2}, \texttt{bootstrap = True}, with best average CV accuracy of
\textbf{96.48\%}.

\clearpage
\section{5-Fold Cross-Validation Performance Comparison}

Table~\ref{tab:fold_compare} reports the fold-wise accuracy of both \emph{tuned} models on the training
set, obtained via \texttt{cross\_val\_score} with the same stratified 5-fold split used during
hyperparameter search. Figure~\ref{fig:fold_plot} visualizes the same comparison.

\begin{table}[h]
\centering
\caption{5-Fold Cross-Validation Accuracy Comparison}
\label{tab:fold_compare}
\begin{tabular}{lcccccc}
\toprule
\textbf{Model} & \textbf{Fold 1} & \textbf{Fold 2} & \textbf{Fold 3} & \textbf{Fold 4} & \textbf{Fold 5} & \textbf{Average} \\
\midrule
Decision Tree  & 0.9451 & 0.9451 & 0.9121 & 0.9231 & 0.9670 & 0.9385 \\
Random Forest  & 0.9670 & 0.9780 & 0.9341 & 0.9560 & 0.9890 & 0.9648 \\
\bottomrule
\end{tabular}
\end{table}

\begin{figure}[h]
  \centering
  \includegraphics[width=0.7\linewidth]{figs/fig_08.png}
  \caption{5-fold CV accuracy per fold — Decision Tree vs.\ Random Forest.}
  \label{fig:fold_plot}
\end{figure}

The Random Forest outperforms the Decision Tree on every single fold and has a higher average accuracy
(96.48\% vs.\ 93.85\%), indicating both better and more \emph{consistent} generalization.

\section{Model Comparison on the Held-out Test Set}

\begin{table}[h]
\centering
\caption{Evaluation Metrics on the Held-out Test Set (Tuned Models)}
\label{tab:eval_metrics}
\begin{tabular}{lcccc}
\toprule
\textbf{Model} & \textbf{Accuracy} & \textbf{Precision} & \textbf{Recall} & \textbf{F1-score} \\
\midrule
Decision Tree (tuned) & 0.9035 & 0.9420 & 0.9028 & 0.9220 \\
Random Forest (tuned) & 0.9561 & 0.9589 & 0.9722 & 0.9655 \\
\bottomrule
\end{tabular}
\end{table}

\begin{figure}[h]
  \centering
  \includegraphics[width=\linewidth]{figs/fig_05.png}
  \caption{Confusion matrices on the held-out test set — Decision Tree (left) vs.\ Random Forest (right).}
  \label{fig:cm}
\end{figure}

\begin{figure}[h]
  \centering
  \includegraphics[width=0.6\linewidth]{figs/fig_06.png}
  \caption{ROC curves — Decision Tree vs.\ Random Forest.}
  \label{fig:roc}
\end{figure}

\begin{figure}[h]
  \centering
  \includegraphics[width=\linewidth]{figs/fig_07.png}
  \caption{Top-10 feature importances — Decision Tree (left) vs.\ Random Forest (right).}
  \label{fig:importance}
\end{figure}

On every metric --- accuracy, precision, recall, F1-score, and ROC-AUC --- the tuned Random Forest
outperforms the tuned Decision Tree on the held-out test set, and its confusion matrix shows fewer
misclassifications of malignant cases (the clinically more costly error type).

\clearpage
\section{Observation Questions}

\textbf{1. How does tree depth affect overfitting in Decision Trees?}
As \texttt{max\_depth} increases, the tree keeps splitting on progressively smaller subsets of the
training data and eventually starts modeling noise rather than genuine patterns. This was directly
observed in the baseline tree, which reached 100\% training accuracy but only 91.23\% test accuracy.
Cross-validation selected a moderate depth (5), which trades a small amount of training accuracy for a
smaller train/test gap and better generalization.

\textbf{2. Which hyperparameter had the greatest impact on performance?}
For the Decision Tree, \texttt{max\_depth} had the largest effect on the accuracy/overfitting
trade-off. For the Random Forest, \texttt{n\_estimators} combined with \texttt{max\_depth} mattered
most --- more trees reduce variance through averaging, while \texttt{max\_features} had a comparatively
smaller marginal effect once enough trees were combined.

\textbf{3. How does Random Forest improve generalization?}
By bagging (training each tree on a bootstrapped sample) and randomly selecting a subset of features at
each split, Random Forest builds many decorrelated trees. Averaging their predictions (majority voting)
cancels out individual trees' variance, producing more stable, better-generalizing predictions than any
single deep tree.

\textbf{4. Did ensemble learning always improve performance? Why or why not?}
Here, yes --- the tuned Random Forest exceeded the tuned Decision Tree on every evaluation metric and
on every cross-validation fold. This is expected because the base learner (Decision Tree) has
comparatively high variance on this dataset; ensembling helps most in exactly this regime. In general,
however, ensembling is not guaranteed to always win: if the base tree is already well-regularized and
close to optimal, or the dataset is very small/simple, the marginal accuracy gain from ensembling can
shrink while computational cost and interpretability get worse.

\section{Conclusion}
Decision Tree and Random Forest models were implemented and evaluated on the Wisconsin Diagnostic
Breast Cancer dataset using 5-fold cross-validation. Hyperparameters were selected based on average
cross-validation performance, ensuring robust generalization rather than overfitting to a single
train/test split. The results demonstrate that Random Forest reduces variance and improves stability
compared to a single Decision Tree: it achieved a higher average cross-validation accuracy (96.48\% vs.\
93.85\%), a higher test-set accuracy (95.61\% vs.\ 90.35\%) and F1-score (0.9655 vs.\ 0.9220), a smaller
train/test accuracy gap, and more consistent fold-wise performance.

\section*{References}
\begin{itemize}
  \item Scikit-learn Documentation: Decision Trees --- \url{https://scikit-learn.org/stable/modules/tree.html}
  \item Scikit-learn Documentation: Random Forest --- \url{https://scikit-learn.org/stable/modules/ensemble.html\#forest}
  \item UCI Machine Learning Repository --- \url{https://archive.ics.uci.edu}
\end{itemize}

\end{document}
